\section{Otvorena pitanja}\label{sec:otvoreno}

\begin{enumerate}[leftmargin=1.6em]
\item \textbf{Koja varijanta modela daje H9?} Hipoteza H9 povezuje $\alpha\neq1$ s indeterminiranim rastom. Minimalni model to ne reproducira: $L_\infty$ ne ovisi o $\alpha$ (P1), a uz ograničen protok (P5) upravo $\alpha<1$ postavlja granicu veličine, dakle suprotno od H9. Ne pomaže ni varijanta u kojoj asimilacija prolazi kroz isto sučelje ($a=\alpha$): uz $a=1$ jedn.~\eqref{eq:mV} u kvazi-ravnoteži daje $dV/dt\propto V$, tj. rast bez asimptote kao kod V1-morfa (T3). Indeterminirani rast tada bi pratio $\alpha=1$, a ne $\alpha\neq1$. Otvoreno je koja pretpostavka može dati H9, npr. rast izvan kvazi-ravnoteže ili drukčije pravilo održavanja ili reprodukcije.
\item \textbf{Homeostaza (H8) i P2.} Prema P2 gustoća spremišta skalira kao $V^{a-\alpha}$. Uz $a=2/3$ i $\alpha=1$ mijenja se deseterostruko između $L=0{,}3$ i $3$~cm, što je teško uskladiti s H8. H8 je zadovoljena samo uz $a\approx\alpha$ (Dodatak~\ref{app:izvodi}). Otvoreno je prolazi li asimilacija doista kroz isto sučelje kao izmjena sa spremištem ili su $a$ i $\alpha$ povezani na drugi način. Odgovor daju mjerenja unosa i spremišta prema veličini (pogl.~\ref{sec:protokol}, Korak 3).
\item \textbf{Regulator koncentracije u transportu (P5).} Uvodi li se, $\alpha$ postaje ograničenje protoka i počinje utjecati na rast.
\item \textbf{κ-DEB kao poseban slučaj tDEB-a.} Uz konstantne $\beta_S$ i $\beta_R$ P1 u kvazi-ravnoteži daje κ-pravilo s $\kappa=\beta_S/(\beta_S+\beta_R)$. Formalni izvod izvan kvazi-ravnoteže ostaje otvoren.
\item \textbf{Dva $\alpha$} (energija i kisik, H5) i njihovo sprezanje u jednadžbama modela.
\item \textbf{$\alpha$ ovisan o stadiju}, tj. veza s akceleracijom tipa M (T3).
\item \textbf{Troškovi transporta.} Minimalni model ne sadrži trošak rada transportne mreže. Otvoreno je koliki je taj trošak, ovisi li o $\alpha$ i koliko pridonosi respiraciji.
\item \textbf{Kvazi-ravnoteža kod mladih.} P1–P3 pretpostavljaju da su $E_T$ i $E_E$ u kvazi-ravnoteži. To je ograda modela: kod mladih, koji relativno brže rastu, spremište zaostaje za ravnotežom (S4), a početni uvjeti važni su za najmanje jedinke (S5). Otvoreno je kako tu ogradu ugraditi u procjenu za vrste s brzim ranim rastom.
\item \textbf{Otvoreni krvotoci.} U pregledu za ovaj priručnik nije nađena teorija analogna WBE-u za otvorene krvotoke, a to je najveća praznina (tablica~\ref{tab:priori}).
\item \textbf{Prazan prostor u literaturi.} Pregledom nije nađen rad koji DEB spaja s razgranatom opskrbnom mrežom ili tretira eksponent površine kao slobodan parametar, osim kroz funkciju korekcije oblika i V1-morfe (T3). To je prilika za objavu.
\end{enumerate}
